% Kompilacja (z katalogu projektu):
%   pdflatex dokumentacja_api.tex
%   pdflatex dokumentacja_api.tex
\documentclass[12pt,a4paper]{report}

\usepackage[utf8]{inputenc}
\usepackage[T1]{fontenc}
\usepackage[polish]{babel}
\usepackage{lmodern}
\usepackage{microtype}
\usepackage{geometry}
\geometry{left=2.5cm,right=2.5cm,top=2.5cm,bottom=2.5cm}

\usepackage{booktabs}
\usepackage{tabularx}
\usepackage{listings}
\usepackage{xcolor}
\usepackage{hyperref}
\usepackage{enumitem}
\usepackage{tikz}
\usepackage{needspace}
\usepackage{float}
\usepackage{etoolbox}
\usetikzlibrary{arrows.meta,positioning,shapes.geometric}

\clubpenalty=10000
\widowpenalty=10000
\displaywidowpenalty=10000
\raggedbottom

\hypersetup{
  colorlinks=true,
  linkcolor=blue!60!black,
  urlcolor=blue!70!black,
  pdftitle={Dokumentacja API Alkozon},
  pdfauthor={Robert Kuzba},
}

\newcommand{\blok}[1]{\needspace{9\baselineskip}#1}
\newcommand{\tabela}[1]{\needspace{11\baselineskip}#1}
\newcommand{\rysunek}[1]{\needspace{14\baselineskip}#1}

\newenvironment{calosc}{%
  \Needspace{14\baselineskip}%
  \begin{samepage}%
}{%
  \end{samepage}%
}
\pretocmd{\section}{\Needspace{14\baselineskip}}{}{}
\pretocmd{\subsection}{\Needspace{11\baselineskip}}{}{}

\lstdefinestyle{shell}{
  basicstyle=\ttfamily\small,
  breaklines=true,
  frame=single,
  backgroundcolor=\color{gray!8},
}
\lstdefinestyle{tree}{
  basicstyle=\ttfamily\footnotesize,
  breaklines=true,
  frame=single,
  backgroundcolor=\color{gray!8},
}

\title{%
  \textbf{API Alkozon}\\[0.4em]
  \large Dokumentacja projektowa --- backend
}
\author{Robert Kuzba}
\date{Legnica, 2026}

\begin{document}

\begin{titlepage}
  \thispagestyle{empty}
  \centering
  \vspace*{1cm}
  {\large Collegium Witelona Uczelnia Państwowa\par}
  \vspace{0.5cm}
  {\large Zaawansowane metody programowania\par}
  \vspace{2.5cm}
  {\LARGE\bfseries API Alkozon\par}
  \vspace{0.8cm}
  {\large Dokumentacja projektowa\par}
  \vspace{0.5cm}
  {\large Temat: Firma produkująca alkohol --- system obsługi zamówień\par}
  \vfill
  {\large Autor:\par}
  {\large Robert Kuzba\par}
  \vspace{1cm}
  {\normalsize Repozytorium: \url{https://github.com/Robertkuzba/api-alcozon}\par}
  \vspace{0.5cm}
  {\normalsize Wdrożenie: \url{https://api-alcozon.onrender.com}\par}
  \vspace{0.3cm}
  {\normalsize Dokumentacja interfejsu: \url{https://api-alcozon.onrender.com/docs}\par}
  \vspace{1.5cm}
  Legnica, 2026
\end{titlepage}

\pagenumbering{roman}
\tableofcontents
\clearpage
\pagenumbering{arabic}

%==============================================================================
\chapter{Wstęp}
%==============================================================================

\blok{%
W ramach projektu zespołowego opracowano system \textbf{Alkozon} --- platformę obsługi sprzedaży produktów alkoholowych. \textbf{API} (\texttt{alcohol-factory-api}) stanowi centralny backend: udostępnia logikę biznesową, przechowuje dane w PostgreSQL oraz łączy klientów \textbf{web} (Next.js), \textbf{mobilny} (React Native) i \textbf{desktop} (Python).

Niniejsze sprawozdanie dotyczy wyłącznie warstwy serwerowej. Zakres obejmuje: założenia funkcjonalne, stos technologiczny, architekturę modułową, bezpieczeństwo, komunikację w czasie rzeczywistym (STOMP), integrację z aplikacjami klienckimi, testy oraz wdrożenie produkcyjne.
}

%==============================================================================
\chapter{Założenia i funkcjonalności}
%==============================================================================

\blok{%
Poniżej zestawiono główne obszary funkcjonalne API (tabela~\ref{tab:funkcje}).
}

\tabela{%
\begin{table}[H]
  \centering
  \caption{Funkcjonalności backendu Alkozon}
  \label{tab:funkcje}
  \begin{tabularx}{\textwidth}{@{}lX@{}}
    \toprule
    \textbf{Obszar} & \textbf{Opis} \\
    \midrule
    Logika centralna & REST \texttt{/api/*} --- wspólny kontrakt dla Web, Mobile i Desktop. \\
    Zamówienia sklepowe & Pełny cykl statusów: \texttt{SUBMITTED} $\rightarrow$ \texttt{IN\_PRODUCTION} $\rightarrow$ \texttt{IN\_PACKING} $\rightarrow$ \texttt{IN\_DELIVERY} $\rightarrow$ \texttt{DELIVERED} (oraz \texttt{CANCELLED}). \\
    Zamówienia własne & Przyjmowanie zapytań z Weba; obsługa przez staff (mobile/desktop). \\
    Dostawy i kurierzy & Rekord \texttt{Delivery} przy \texttt{IN\_DELIVERY}; przypisanie kuriera (MANAGER); lista zamówień kuriera. \\
    Real-time & WebSocket STOMP --- natychmiastowe powiadomienia o zmianie statusu (Web, magazyn, kurier). \\
    Push (mobilka) & Firebase Cloud Messaging po przypisaniu dostawy i zmianach statusu. \\
    Moduł HR & Rejestracja czasu pracy (clock-in/out, przerwa) --- mobile; agregacja godzin --- desktop (MANAGER). \\
    Raporty & Sprzedaż, stan magazynu, podsumowanie czasu pracy --- desktop (MANAGER). \\
    Magazyn / produkty & Katalog, stany, uzupełnienia, administracja użytkowników i ogłoszeń. \\
    Dokumentacja API & OpenAPI 3 + Swagger UI pod adresem \texttt{/docs}. \\
    \bottomrule
  \end{tabularx}
\end{table}
}

%==============================================================================
\chapter{Stack technologiczny}
%==============================================================================

\tabela{%
\begin{table}[H]
  \centering
  \caption{Stos technologiczny API}
  \label{tab:stack}
  \begin{tabularx}{\textwidth}{@{}lX@{}}
    \toprule
    \textbf{Warstwa} & \textbf{Technologia} \\
    \midrule
    Język & Java 21 LTS \\
    Framework & Spring Boot 4 (Web, Security, Data JPA, WebSocket/STOMP) \\
    Baza danych & PostgreSQL; migracje Flyway \\
    Bezpieczeństwo & Spring Security, JWT (access + refresh), BCrypt \\
    Rate limiting & Bucket4j (filtr servletowy) \\
    Dokumentacja & springdoc-openapi (Swagger UI) \\
    Powiadomienia push & Firebase Admin SDK (FCM) \\
    Konteneryzacja & Docker (\texttt{Dockerfile}, \texttt{docker-compose} dla Postgresa lokalnie) \\
    CI & GitHub Actions --- \texttt{mvn verify} \\
    \bottomrule
  \end{tabularx}
\end{table}
}

%==============================================================================
\chapter{Architektura}
%==============================================================================

\section{Struktura modułów}

\blok{%
Kod źródłowy podzielono na moduły domenowe w pakiecie \texttt{com.alcoholfactory.api.modules.*}:
}

\begin{lstlisting}[style=tree]
alcohol-factory-api/
  src/main/java/.../modules/
    auth/        # logowanie, JWT, 2FA staff, reset hasla
    order/       # zamowienia sklepowe i custom
    delivery/    # dostawy, kurierzy
    product/     # katalog
    inventory/   # stany magazynowe
    warehouse/   # uzupelnienia (MANAGER)
    hr/          # work-log
    report/      # raporty (MANAGER)
    admin/       # uzytkownicy, oferty, ogloszenia
    device/      # tokeny FCM
    notification/# STOMP + FCM
  src/main/resources/db/migration/  # Flyway V1..Vn
\end{lstlisting}

\section{Przepływ w systemie rozproszonym}

\rysunek{%
\begin{figure}[H]
  \centering
  \begin{tikzpicture}[
    node distance=0.9cm and 1.1cm,
    box/.style={draw, rounded corners, minimum width=2.2cm, minimum height=0.65cm, align=center, font=\small},
    api/.style={draw, rounded corners, minimum width=2.6cm, minimum height=0.8cm, align=center, font=\small\bfseries, fill=blue!6},
    db/.style={cylinder, draw, shape border rotate=90, aspect=0.25, minimum height=0.85cm, align=center, font=\small},
  ]
    \node[box] (web) {Web};
    \node[box, right=of web] (mob) {Mobile};
    \node[box, right=of mob] (desk) {Desktop};
    \node[api, below=1.4cm of mob] (api) {API REST + STOMP};
    \node[db, below=of api] (db) {PostgreSQL};
    \draw[-{Stealth[length=2mm]}] (web) -- (api);
    \draw[-{Stealth[length=2mm]}] (mob) -- (api);
    \draw[-{Stealth[length=2mm]}] (desk) -- (api);
    \draw[-{Stealth[length=2mm]}] (api) -- (db);
  \end{tikzpicture}
  \caption{API jako centralny hub komunikacji}
  \label{fig:arch}
\end{figure}
}

REST obsługuje operacje synchroniczne (złożenie zamówienia, zmiana statusu, raporty). STOMP uzupełnia to zdarzeniami push do subskrybentów (np.\ klient web po zalogowaniu, magazyn po temacie \texttt{/topic/orders/staff}, kurier po \texttt{/user/queue/courier-deliveries}).

%==============================================================================
\chapter{Realizacja funkcjonalności}
%==============================================================================

\begin{calosc}
\section{Zamówienia i synchronizacja statusów}

Zamówienia sklepowe (\texttt{/api/orders}) oraz zamówienia własne (\texttt{/api/custom-orders}) korzystają ze wspólnego zestawu statusów (\texttt{OrderStatus}). Zmiana statusu przez pracownika (\texttt{PATCH .../status}) aktualizuje rekord w bazie, a następnie uruchamia \textbf{notifier} czasu rzeczywistego (\texttt{OrderRealtimeNotifier} / \texttt{CustomOrderRealtimeNotifier}), który publikuje zdarzenie STOMP i --- w razie potrzeby --- powiadomienie FCM.

Przy przejściu do \texttt{IN\_DELIVERY} tworzony jest rekord \texttt{Delivery}. Manager przypisuje kuriera endpointem \texttt{PATCH /api/deliveries/\{id\}/assign} (tylko rola \texttt{MANAGER}). Kurier pobiera swoje zlecenia przez \texttt{GET /api/orders/for-courier/\{courierUserId\}\}.
\end{calosc}

\begin{calosc}
\section{Powiadomienia w czasie rzeczywistym}

Broker STOMP jest wystawiony pod \texttt{/ws} (protokół \texttt{wss} na produkcji). Przy połączeniu klient przekazuje token JWT w nagłówku \texttt{Authorization} komendy \texttt{CONNECT}. Główne destynacje:
\begin{itemize}[noitemsep]
  \item \texttt{/topic/orders/staff} --- magazyn / staff (zmiany statusów),
  \item \texttt{/topic/orders/dispatch} --- desktop managera (np.\ nowa dostawa bez kuriera),
  \item \texttt{/user/queue/order-updates} --- klient \texttt{CUSTOMER} (śledzenie własnego zamówienia),
  \item \texttt{/user/queue/courier-deliveries} --- kurier po przypisaniu.
\end{itemize}
\end{calosc}

\begin{calosc}
\section{Moduł HR (czas pracy)}

Dla roli \texttt{EMPLOYEE} (aplikacja mobilna):
\begin{itemize}[noitemsep]
  \item \texttt{POST /api/work-log/clock-in} --- rozpoczęcie pracy,
  \item \texttt{POST /api/work-log/clock-out} --- zakończenie,
  \item \texttt{POST /api/work-log/break/start}, \texttt{break/end} --- przerwa,
  \item \texttt{GET /api/work-log/my?from=\&to=} --- historia własna.
\end{itemize}
Manager (desktop) korzysta z \texttt{GET /api/reports/employees/work-summary} oraz \texttt{GET /api/work-log/reports/summary}.
\end{calosc}

\begin{calosc}
\section{Analityka i raporty}

Endpointy pod \texttt{/api/reports} (rola \texttt{MANAGER}):
\begin{itemize}[noitemsep]
  \item \texttt{GET /reports/sales} --- suma sprzedaży w przedziale dat,
  \item \texttt{GET /reports/inventory} --- snapshot stanów magazynowych,
  \item \texttt{GET /reports/employees/work-summary} --- agregacja godzin pracy.
\end{itemize}
Dane zasilają wykresy w aplikacji desktopowej.
\end{calosc}

\begin{calosc}
\section{Dokumentacja interfejsu (OpenAPI)}

Interaktywna dokumentacja dostępna jest bez dodatkowej instalacji:
\begin{itemize}[noitemsep]
  \item lokalnie: \url{http://localhost:8080/docs},
  \item produkcja: \url{https://api-alcozon.onrender.com/docs}.
\end{itemize}
Kontrakt REST jest źródłem prawdy dla zespołów frontendowych (generowanie typów, testy integracyjne w Swaggerze).
\end{calosc}

%==============================================================================
\chapter{Bezpieczeństwo}
%==============================================================================

\blok{%
Bezpieczeństwo API opiera się na uwierzytelnianiu tokenowym, kontroli ról oraz walidacji i limitach po stronie serwera (tabela~\ref{tab:security}).
}

\tabela{%
\begin{table}[H]
  \centering
  \caption{Mechanizmy bezpieczeństwa API}
  \label{tab:security}
  \begin{tabularx}{\textwidth}{@{}lX@{}}
    \toprule
    \textbf{Mechanizm} & \textbf{Realizacja} \\
    \midrule
    JWT & Access token (krótki TTL) + refresh token; weryfikacja w filtrze Spring Security. \\
    Role (RBAC) & \texttt{GUEST}, \texttt{CUSTOMER}, \texttt{EMPLOYEE}, \texttt{MANAGER} --- adnotacje \texttt{@PreAuthorize} na kontrolerach. \\
    2FA staff & Logowanie pracowników: kod 4-cyfrowy na nowe \texttt{deviceId}; zaufane urządzenia w tabeli \texttt{trusted\_devices}. \\
    Walidacja wejścia & Jakarta Validation na DTO żądań --- druga linia obrony względem klientów. \\
    Hasła & BCrypt (\texttt{PasswordEncoder}); reset hasła pracownika --- losowe hasło tymczasowe e-mailem. \\
    Rate limiting & Bucket4j: logowanie 10/15\,min/IP; wyszukiwarka produktów 5\,req/s/IP; śledzenie zamówienia 30/min/IP. \\
    CORS & Whitelist originów (\texttt{app.cors.allowed-origins}) --- m.in.\ Vercel i localhost. \\
    WebSocket & Autoryzacja przy \texttt{CONNECT} i \texttt{SUBSCRIBE} --- token JWT w nagłówku STOMP. \\
    \bottomrule
  \end{tabularx}
\end{table}
}

\begin{calosc}
\section{Podział ról (przykłady)}

\begin{itemize}[noitemsep]
  \item \texttt{CUSTOMER} --- składanie zamówień, śledzenie, konto.
  \item \texttt{EMPLOYEE} --- zmiana statusów, magazyn, własny work-log, dostawy przypisane do kuriera.
  \item \texttt{MANAGER} --- wszystko co EMPLOYEE oraz: przypisanie kuriera, raporty, administracja użytkowników/produktów, magazyn zaawansowany.
\end{itemize}
\end{calosc}

%==============================================================================
\chapter{Integracja z klientami i wdrożenie}
%==============================================================================

\begin{calosc}
\section{Kontrakt REST}

Wszystkie endpointy biznesowe mają prefiks \texttt{/api}. Przykłady:
\begin{itemize}[noitemsep]
  \item \texttt{POST /api/auth/login} --- klient (CUSTOMER),
  \item \texttt{POST /api/auth/staff/login} --- pracownik (2FA),
  \item \texttt{POST /api/orders}, \texttt{PATCH /api/orders/\{id\}/status},
  \item \texttt{POST /api/custom-orders}, \texttt{PATCH /api/custom-orders/\{id\}/status}.
\end{itemize}
\end{calosc}

\begin{calosc}
\section{Uruchomienie lokalne}

\begin{lstlisting}[style=shell]
docker compose up -d
./mvnw spring-boot:run
\end{lstlisting}
Swagger: \url{http://localhost:8080/docs}. WebSocket: \texttt{ws://localhost:8080/ws}.

Po pierwszym starcie (profil nie-test) tworzone są konta demo (m.in.\ \texttt{manager@example.com}, \texttt{employee@example.com}) oraz przykładowy produkt --- szczegóły w \texttt{README.md}.
\end{calosc}

\begin{calosc}
\section{Wdrożenie produkcyjne}

API wdrożone jest na platformie \textbf{Render} (\url{https://api-alcozon.onrender.com}), baza danych --- \textbf{Neon PostgreSQL}. Migracje Flyway uruchamiane przy starcie aplikacji. Zmienne środowiskowe obejmują m.in.: \texttt{DATABASE\_URL}, \texttt{JWT\_SECRET}, \texttt{APP\_CORS\_ALLOWED\_ORIGINS}, konfigurację SMTP (2FA, reset hasła) oraz opcjonalnie \texttt{FIREBASE\_SERVICE\_ACCOUNT\_JSON} dla pushy mobilnych.

Build obrazu Docker:
\begin{lstlisting}[style=shell]
./mvnw -q package -DskipTests
docker build -t alcohol-api .
\end{lstlisting}
\end{calosc}

%==============================================================================
\chapter{Testowanie}
%==============================================================================

\blok{%
Testy automatyczne uruchamiane poleceniem \texttt{./mvnw verify} (wymaga Dockera dla Testcontainers w testach integracyjnych). Pakiet \texttt{src/test/java} obejmuje m.in.:
}

\begin{itemize}[noitemsep]
  \item autoryzację i 2FA staff (\texttt{StaffAuthMvcTest}),
  \item rate limiting logowania (\texttt{AuthLoginRateLimitMvcTest}),
  \item przejścia statusów zamówień i dostaw (\texttt{OrderStatusTransitionMvcTest}, \texttt{DeliveryAssignAuthorizationMvcTest}),
  \item zamówienia własne i notyfikacje (\texttt{CustomOrderStatusMvcTest}, \texttt{DefaultOrderRealtimeNotifierTest}),
  \item reset hasła pracownika (\texttt{EmployeePasswordResetMvcTest}).
\end{itemize}

CI w repozytorium GitHub uruchamia pełną weryfikację Maven przy każdym pushu na gałąź główną.

%==============================================================================
\chapter{Podsumowanie}
%==============================================================================

\blok{%
Opracowano \textbf{centralne API Alkozon} w technologii Java~21 / Spring Boot~4, integrujące aplikację webową, mobilną i desktopową. Backend realizuje pełny cykl zamówień (sklep i konfigurator własny), moduł dostaw z kurierami, powiadomienia STOMP i FCM, rejestrację czasu pracy oraz raporty dla managera. Zabezpieczenia obejmują JWT z rolami, walidację DTO, BCrypt, rate limiting (Bucket4j) oraz dwuskładnikowe logowanie staff.

API jest udokumentowane w OpenAPI (\texttt{/docs}) i wdrożone pod adresem \url{https://api-alcozon.onrender.com}. Dalszy rozwój obejmuje m.in.\ rozszerzenie raportów, utwardzenie polityki 2FA na środowiskach deweloperskich oraz dalszą automatyzację testów kontraktu z klientami.
}

\end{document}
